% !TEX program = lualatex
\documentclass[11pt,a4paper]{article}

% ========= Пакеты =========
\usepackage[english, bidi=default, layout=counters]{babel}
\usepackage{amsmath,amssymb,amsfonts,amsthm,mathtools}
\usepackage{geometry}
\geometry{left=1.25cm,right=1.25cm,top=1.25cm,bottom=3.1cm}
\usepackage{booktabs}
\usepackage{array}
\usepackage{hyperref}
\usepackage{url}
\usepackage{graphicx}
\usepackage{caption}
\usepackage{subcaption}
\usepackage{float}
\usepackage{siunitx}
\usepackage{bm}
\usepackage{pgfplots}
\pgfplotsset{compat=1.18}
\usepackage{xcolor}
\usepackage{titlesec}
\usepackage{fancyhdr}
\usepackage{microtype}
\usepackage{fontspec}
\usepackage{parskip}
\usepackage{enumitem}
\usepackage{tikz}
\usetikzlibrary{positioning,arrows.meta}
% ========= Числовой стиль natbib =========
\usepackage[numbers]{natbib}
\usepackage{xurl}

% ========= Язык (русский как основной) =========
\babelprovide[import, main]{russian}
\babelprovide[import, onchar=ids fonts]{english}

% ========= Шрифты =========
\defaultfontfeatures{Ligatures=TeX}
\babelfont[russian]{rm}[Renderer=Harfbuzz]{Times New Roman}
\babelfont[russian]{sf}[Renderer=Harfbuzz]{Arial}
\babelfont[russian]{tt}[Renderer=Harfbuzz]{Courier New}
\babelfont[english]{rm}{Times New Roman}
\babelfont[english]{sf}{Arial}
\babelfont[english]{tt}{Courier New}

% ========= Цвета =========
\definecolor{skyblue}{RGB}{0,102,204}
\definecolor{darkblue}{RGB}{0,0,139}
\definecolor{gold}{RGB}{255,215,0}

% ========= Макросы YUCT =========
\newcommand{\Keff}{K_{\mathrm{eff}}}
\newcommand{\Dir}{\mathcal{D}}
\newcommand{\Mpl}{M_{\mathrm{Pl}}}
\newcommand{\Lpl}{l_{\mathrm{Pl}}}
\newcommand{\tP}{t_{\mathrm{P}}}
\newcommand{\Sodd}{S_{\mathrm{odd}}}
\newcommand{\Seven}{S_{\mathrm{even}}}
\newcommand{\kappac}{\kappa_{c}}
\newcommand{\betaf}{\beta}
\newcommand{\lagr}{\mathcal{L}}
\newcommand{\dint}{D\!+\!I\!\cdot\!R}
\newcommand{\CCU}{\text{КСЕ}}

% ========= Оформление заголовков =========
\titleformat{\section}{\LARGE\bfseries\color{darkblue}}{\thesection}{1em}{}
\titleformat{\subsection}{\normalsize\bfseries\color{darkblue}}{\thesubsection}{1em}{}
\titleformat{\subsubsection}{\normalsize\itshape}{\thesubsubsection}{1em}{}

% ========= Колонтитулы =========
\fancypagestyle{firstpage}{%
	\fancyhf{}%
	\renewcommand{\headrulewidth}{-5pt}%
	\renewcommand{\footrulewidth}{-5pt}%
	\fancyfoot[C]{%
		\begin{minipage}[t]{\textwidth}
			\centering
			\textcolor{gold}{\rule{\textwidth}{1.6pt}}\\[5pt]
			{\small \textsuperscript{1}Yakushev Research, \url{https://yuct.org/}} \hfill
			{\small Протокол YPSDC: \url{https://ypsdc.com/}}\\[-2pt]
			\textcolor{gold}{\rule{\textwidth}{1.6pt}}\\[5pt]
			\textbf{Объединённая теория координации Якушева 2026} \hfill \thepage\\[10pt]
		\end{minipage}%
	}%
}

\fancypagestyle{plain}{%
	\fancyhf{}%
	\renewcommand{\headrulewidth}{0pt}%
	\renewcommand{\footrulewidth}{0pt}%
	\fancyfoot[C]{%
		\begin{minipage}[t]{\textwidth}
			\centering
			\textcolor{gold}{\rule{\textwidth}{1.6pt}}\\[4pt]
			\textbf{Объединённая теория координации Якушева 2026} \hfill \thepage\\[4pt]
		\end{minipage}%
	}%
}

\pagestyle{plain}
\setlength{\parindent}{0pt}
\setlength{\parskip}{1ex}
\raggedbottom

\hypersetup{
	colorlinks=true,
	linkcolor=blue,
	citecolor=blue,
	urlcolor=blue,
}

% ========= Шапка (логотип YUCT) =========
\newcommand{\makeheader}{%
	\noindent
	\begin{minipage}{0.17\textwidth}
		\raggedright
		{\sffamily\bfseries\color{skyblue}\fontsize{46}{46}\selectfont YUCT}%
	\end{minipage}%
	\hspace{30pt}
	\begin{minipage}{0.32\textwidth}
		\raggedright
		{\sffamily\fontsize{11}{13}\selectfont Объединённая\\ теория\\ координации Якушева}%
	\end{minipage}%
	\hfill
	\begin{minipage}{0.38\textwidth}
		\raggedleft
		{\sffamily\bfseries Гипотеза / Дорожная карта}\\ 
		{\sffamily Версия 2.0 / Июнь 2026}\\ 
	\end{minipage}
	\par\vspace{5pt}
	\noindent\textcolor{gold}{\rule{\textwidth}{1.6pt}}
	\par\vspace{0pt}
}

\begin{document}
	\thispagestyle{firstpage}
	\makeheader
	
	\vspace{0cm}
	{\LARGE\bfseries Сильная CP-проблема:\\[0.1cm] конкретная алгебраическая гипотеза в рамках YUCT\\[0.2cm] \Large Дорожная карта теоретической и экспериментальной проверки \par}
	\vspace{0.2cm}
	
	{\large Алексей В. Якушев\textsuperscript{1}} \\
	\vspace{0.0cm}
	
	{\color{darkblue}\textbf{\LARGE Аннотация}} \par
	{\large
		\noindent
		Сильная CP-проблема — почему нарушающий CP параметр \(\theta\) в КХД меньше \(10^{-10}\) — является последней неразгаданной загадкой Стандартной модели в Объединённой теории координации Якушева (YUCT). В настоящей работе мы предлагаем конкретную алгебраическую гипотезу, которая естественным образом даёт необходимое подавление. Мы предполагаем, что эффективная глубина координации параметра \(\theta\) равна
		\[
		L_\theta = \frac{1}{2}\,L_W \;+\; \frac{S_{\mathrm{odd}}}{S_{\mathrm{even}}},
		\]
		где \(L_W = 96\) — глубина \(W\)-бозона, а \(S_{\mathrm{odd}}/S_{\mathrm{even}} = 3/2\) — фундаментальная константа теории. Эта формула не содержит свободных параметров и даёт \(L_\theta = 49.5\), что соответствует \(\theta \sim (2/3)^{49.5} \approx 2.3 \times 10^{-10}\). Мы обсуждаем физическую мотивировку данного выражения, встраиваем его в более широкий алгебраический каркас YUCT и намечаем план экспериментальной проверки — прежде всего через электрический дипольный момент нейтрона (ЭДМН), — способный фальсифицировать гипотезу. Предложение формулируется как чёткая, математически точная гипотеза, ожидающая подтверждения или опровержения будущими данными.
	}
	
	\vspace{0.1cm}
	\noindent
	{\color{darkblue}\textbf{Ключевые слова:}} YUCT, сильная CP-проблема, алгебраическая гипотеза, глубина координации, электрический дипольный момент нейтрона, открытый исследовательский план.
	\vspace{0.4cm}
	
	\clearpage
	\tableofcontents
	\newpage
	
	\section{Введение}
	\label{sec:intro}
	
	Объединённая теория координации Якушева (YUCT) уже дала количественные, не содержащие свободных параметров решения для проблемы иерархий, проблемы космологической постоянной, структуры фермионных масс и ряда других крупных загадок фундаментальной физики~\cite{Y1,Y3,Y4}. Единственным исключением оставалась сильная CP-проблема: неестественно малая величина нарушающего CP параметра \(\theta\) в квантовой хромодинамике (КХД), экспериментальный верхний предел для которой составляет \(|\theta| \lesssim 10^{-10}\)~\cite{nEDM2020}.
	
	В рамках YUCT ранние попытки приписать \(\theta\) глубину координации с помощью специально вводимого целого числа \(k_\theta\) были неудовлетворительными, так как сводились к подгонке, а не к структурному предсказанию. В настоящей заметке мы предлагаем иной подход: мы используем взаимодействие между упорядоченным (фермионным) и хаотическим (топологическим) секторами вакуума КХД — взаимодействие, которое YUCT способна описывать количественно благодаря своему единому алгебраическому описанию порядка и хаоса~\cite{Y2,Feig}.
	
	Мы не утверждаем, что даём строгий вывод из 19-мерного лагранжиана; такой вывод является долгосрочной целью. Вместо этого мы формулируем точную алгебраическую гипотезу, которая естественно вытекает из уже присутствующих в YUCT чисел, не содержит подгоночных параметров и даёт фальсифицируемые предсказания для электрического дипольного момента нейтрона (ЭДМН) и массы аксиона. Данная гипотеза представляет собой чётко определённую цель для будущих теоретических и экспериментальных исследований.
	
	\section{Константы порядка и хаоса в YUCT}
	\label{sec:oc}
	
	\subsection{Порядок}
	Универсальный закон ошибки
	\[
	\varepsilon = \kappa_c\,\alpha\,K_{\mathrm{eff}}^{-\beta},\qquad \beta = \frac{2}{3},\quad \kappa_c = \frac{1}{3},
	\]
	вместе с иерархической структурой координации приводит к замкнутому алгебраическому циклу
	\[
	S_{\mathrm{odd}} = \frac{6}{5} = 1.2,\quad S_{\mathrm{even}} = \frac{4}{5} = 0.8,\quad 
	\frac{S_{\mathrm{odd}}}{S_{\mathrm{even}}} = \frac{3}{2} = \frac{1}{\beta}.
	\]
	Все физические массы выражаются через планковскую массу и глубину координации \(L\):
	\[
	m \approx M_{\mathrm{Pl}}\left(\frac{2}{3}\right)^L.
	\]
	В частности, масса \(W\)-бозона соответствует \(L_W = 96\), т.е. полному числу вейлевских фермионных степеней свободы Стандартной модели (включая правые нейтрино)~\cite{Y3}.
	
	\subsection{Хаос}
	Универсальные константы Фейгенбаума
	\[
	\delta \approx 4.6692,\qquad \alpha \approx 2.5029,
	\]
	описывают каскад удвоений периода на пути к хаосу. В YUCT они алгебраически связаны с константами порядка через золотое сечение \(\varphi = (1+\sqrt{5})/2\) и число \(\pi\):
	\[
	2\alpha^2 \approx \pi\delta,\qquad
	\pi \approx S_{\mathrm{odd}}\varphi^2 \quad \Longrightarrow \quad 
	2\alpha^2 \approx (S_{\mathrm{odd}}\varphi^2)\,\delta .
	\]
	Эта связь показывает, что константы Фейгенбаума не являются чужеродными для YUCT; они принадлежат той же расширенной алгебраической системе~\cite{Y2}.
	
	\subsection{Квантовый вакуум как сопряжение порядка и хаоса}
	Вакуум КХД представляет собой арену сосуществования упорядоченного аспекта (кварковый конденсат, нарушение киральной симметрии) и хаотического аспекта (инстантонные флуктуации, топологическая восприимчивость). Параметр \(\theta\) связан с плотностью топологического заряда \(G\widetilde{G}\), которая получает вклады от обоих секторов. Поэтому естественно ожидать, что эффективная глубина \(L_\theta\) является функцией глубин, характеризующих упорядоченную и хаотическую составляющие.
	
	В предыдущих работах мы вводили предварительные глубины \(L_{\mathrm{ord}}\) (фермионный сектор) и \(L_{\mathrm{ch}}\) (топологический/хаотический сектор), однако их простая линейная комбинация не воспроизводила требуемое подавление \(|\theta| \sim 10^{-10}\). Это побудило нас искать более структурную комбинацию.
	
	\section{Алгебраическая гипотеза о \(L_\theta\)}
	\label{sec:hypothesis}
	
	\subsection{Формулировка гипотезы}
	
	Мы выдвигаем гипотезу, что эффективная глубина координации параметра \(\theta\) даётся простой формулой
	\begin{equation}
		\boxed{L_\theta = \frac{1}{2}\,L_W \;+\; \frac{S_{\mathrm{odd}}}{S_{\mathrm{even}}}},
		\label{eq:Ltheta}
	\end{equation}
	где
	\begin{itemize}
		\item \(L_W = 96\) — глубина \(W\)-бозона (базовый масштаб электрослабой шкалы);
		\item \(S_{\mathrm{odd}}/S_{\mathrm{even}} = 3/2 = 1.5\) — универсальное отношение, управляющее межпоколенческими массовыми факторами и возникающее в первичном алгебраическом цикле.
	\end{itemize}
	
	Подставляя эти числа, получаем
	\[
	L_\theta = \frac{96}{2} + 1.5 = 48 + 1.5 = 49.5.
	\]
	
	Соответствующее значение параметра \(\theta\), с точностью до множителя порядка единицы, возникающего при инстантонных вычислениях, составляет
	\[
	\theta \sim \left(\frac{2}{3}\right)^{L_\theta} = \left(\frac{2}{3}\right)^{49.5} \approx 2.3 \times 10^{-10}.
	\]
	Эта величина попадает точно в экспериментально разрешённое окно \(|\theta| \lesssim 10^{-10}\)~\cite{nEDM2020}.
	
	\subsection{Физическая интерпретация}
	
	Почему возникает именно такая комбинация?
	\begin{enumerate}
		\item \textbf{Множитель \(1/2\)}: электрослабый масштаб задаётся вакуумным средним \(v \approx 246\) ГэВ, которое получается из фундаментальной глубины 96 через геометрический фактор \(\xi_W \approx 0.53\). Глубина 48 соответствует масштабу \(\sqrt{v M_{\mathrm{Pl}}}\) — среднему геометрическому слабого и планковского масштабов, — то есть масштабу, на котором, как ожидается, пересекаются вклады фермионов и бозонов в энергию вакуума. В YUCT известно, что этот масштаб играет особую роль в координационной динамике хиггсовского конденсата (Приложение \(\Lambda\)).
		\item \textbf{Сдвиг \(S_{\mathrm{odd}}/S_{\mathrm{even}} = 3/2\)}: отношение \(3/2 = 1/\beta\) количественно характеризует неустранимую асимметрию между упорядоченными (нечётные слои) и неупорядоченными (чётные слои) корреляциями в любой иерархически координированной системе. Его появление в \(L_\theta\) указывает на то, что CP-нарушающий член возникает из малого, неизбежного дисбаланса между секторами вакуума, в которых доминирует порядок, и секторами, в которых доминирует хаос. Этот сдвиг является аддитивным, а не мультипликативным, что отражает дискретную, слоистую структуру координационной сети.
	\end{enumerate}
	
	Таким образом, уравнение~(\ref{eq:Ltheta}) можно прочитать так: глубина \(\theta\) равна половине слабой глубины плюс универсальный сдвиг «порядок–хаос». Оно не содержит свободных параметров и использует лишь числа, уже независимо определённые в YUCT.
	
	\subsection{Непосредственные следствия}
	
	\paragraph{Электрический дипольный момент нейтрона.}
	ЭДМН связан с \(\theta\) соотношением \(d_n \approx 3.6 \times 10^{-16}\,\theta\; e\!\cdot\!\mathrm{см}\) (с точностью до адронных неопределённостей)~\cite{nEDM2020}. При \(\theta \approx 2.3 \times 10^{-10}\) получаем
	\[
	d_n \approx 8.3 \times 10^{-26}\; e\!\cdot\!\mathrm{см}.
	\]
	Текущий экспериментальный предел составляет \(d_n < 1.8 \times 10^{-26}\; e\!\cdot\!\mathrm{см}\) (90\% доверительный уровень). Эксперименты следующего поколения PSI, SNS и TUCAN нацелены на чувствительность порядка \(10^{-28}\; e\!\cdot\!\mathrm{см}\). Нулевой результат на этом уровне исключит гипотезу; положительное обнаружение стало бы драматическим подтверждением.
	
	\paragraph{Физика аксиона.}
	Если сильная CP-проблема решается механизмом Печчеи–Квинн, масса аксиона \(m_a\) связана с \(\theta\) и постоянной распада аксиона \(f_a\). Уравнение~(\ref{eq:Ltheta}) фиксирует \(f_a\) (а следовательно, и \(m_a\)) через ту же логику глубины координации. Это даёт дополнительные экспериментальные цели для ADMX и других поисков аксиона.
	
	\section{Дорожная карта проверки}
	\label{sec:roadmap}
	
	Гипотеза математически точна и фальсифицируема. Мы предлагаем следующий план:
	
	\subsection{Теоретические задачи}
	\begin{enumerate}
		\item \textbf{Вывод из 19-мерного лагранжиана.}
		Необходимо осуществить редукцию эффективного действия 19-мерной материнской теории YUCT к стандартному члену КХД \(G\widetilde{G}\). Цель — показать, что при правильной идентификации вакуумного состояния коэффициент в редукции естественно даёт глубину \(L_\theta = 49.5\). Это сложная, но чётко поставленная математическая задача.
		\item \textbf{Решёточные калибровочные теории.}
		Существующие решёточные вычисления КХД позволяют рассчитывать топологическую восприимчивость и зависимость энергии вакуума от \(\theta\). Анализ, вдохновлённый YUCT, должен проверить, отражается ли алгебраическая структура уравнения~(\ref{eq:Ltheta}) в скейлинговом поведении решёточных результатов при изменении константы связи или в распределении кластеров топологического заряда.
		\item \textbf{Уточнение предсказания для ЭДМН.}
		Адронные матричные элементы, связывающие \(\theta\) и \(d_n\), известны из киральной теории возмущений и решёточной КХД. Специализированные вычисления этих матричных элементов в сочетании со значением \(\theta\), даваемым YUCT, позволят получить острое предсказание, которое можно будет напрямую сравнивать с ограничениями экспериментов следующего поколения по ЭДМН.
	\end{enumerate}
	
	\subsection{Экспериментальные задачи}
	\begin{enumerate}
		\item \textbf{Эксперименты по ЭДМН (PSI, SNS, TUCAN).}
		Измеренный ЭДМН выше \(10^{-26}\;e\!\cdot\!\mathrm{см}\) будет согласовываться с гипотезой; предел ниже \(10^{-28}\;e\!\cdot\!\mathrm{см}\) исключит её.
		\item \textbf{Поиски аксиона (ADMX, MADMAX, IAXO).}
		Если реализуется механизм Печчеи–Квинн, то после оценки \(f_a\) из \(L_\theta\) масса аксиона, предсказываемая YUCT, должна быть вычислимой. Специализированные поиски в соответствующем массовом окне позволят проверить гипотезу.
		\item \textbf{Космологические отпечатки.}
		Если интерференция порядка и хаоса фундаментальна, она может также влиять на фазовые переходы в ранней Вселенной, оставляя специфические отпечатки в реликтовом излучении или в распределении тёмной материи. Эти эффекты можно ограничить с помощью данных Planck, Euclid и будущего эксперимента CMB-S4.
	\end{enumerate}
	
	\section{Приглашение к сотрудничеству}
	\label{sec:invite}
	
	Полный каркас YUCT, включая полный лагранжиан, все приложения и обширную базу данных глубин координации, находится в открытом доступе на Zenodo~\cite{Y1,Y2}. Мы приглашаем физиков, математиков и вычислительных специалистов принять участие в выводе, проверке или фальсификации предложенной здесь гипотезы. Автор (Алексей Якушев, \href{mailto:alexey@yakushev.eu}{alexey@yakushev.eu}) приветствует корреспонденцию, предложения и проекты сотрудничества.
	
	\section{Заключение}
	\label{sec:conclusion}
	
	Мы предложили конкретную, не содержащую свободных параметров алгебраическую формулу для эффективной глубины параметра \(\theta\) КХД в рамках Объединённой теории координации Якушева:
	\[
	L_\theta = \frac{1}{2}L_W + \frac{S_{\mathrm{odd}}}{S_{\mathrm{even}}} = 49.5,
	\]
	дающую \(\theta \approx 2.3 \times 10^{-10}\). Эта гипотеза математически проста, физически мотивирована и допускает прямую фальсификацию в готовящихся экспериментах по ЭДМН и поиску аксионов. Она заменяет ранние попытки, опиравшиеся на произвольное назначение целых чисел, и открывает ясный путь к полному решению сильной CP-проблемы внутри YUCT. Выдержит ли гипотеза экспериментальную проверку, покажет будущее; её ценность состоит в том, что она даёт точную, проверяемую догадку, направляющую дальнейшие исследования.
	
	\vspace{1cm}
	\noindent\textbf{Дата:} Июнь 2026 г.\\
	\textbf{Версия:} 2.0\\
	\textbf{Статус:} Гипотеза с конкретным алгебраическим предсказанием; открыта для экспериментальной проверки и теоретического вывода.
	
	\newpage
	\begin{thebibliography}{99}
		\bibitem{Y1} A.V. Yakushev, \textit{Yakushev Unified Coordination Theory (YUCT)}, Zenodo, 2026. DOI:10.5281/zenodo.18444599.
		\bibitem{Y2} A.V. Yakushev, Appendix AF: The Algebraic Framework of Reality in YUCT, 2026.
		\bibitem{Y3} A.V. Yakushev, Appendix \(\Lambda\): Resolution of the Hierarchy and Cosmological Constant Problems, 2026.
		\bibitem{Y4} A.V. Yakushev, Appendix J: Complete Derivation of Standard Model Flavor Parameters, 2026.
		\bibitem{Feig} M.J. Feigenbaum, \textit{J. Stat. Phys.} \textbf{19}, 25 (1978).
		\bibitem{nEDM2020} C. Abel et al. (nEDM Collaboration), \textit{Phys. Rev. Lett.} \textbf{124}, 081803 (2020).
	\end{thebibliography}
	
\end{document}